\chapter{Configurazione: QGroundControl e firmware PX4}
\label{cap:configurazione}

Una volta assemblato l'hardware, il velivolo deve essere reso \emph{volabile}:
occorre dire al firmware com'\`e fatto il drone, tarare i sensori e abilitare la
registrazione dei dati. Tutte queste operazioni sono state svolte tramite la
\textbf{ground station QGroundControl} (QGC), l'applicazione desktop che dialoga
con il firmware \textbf{PX4} del Pixhawk 6X: permette di leggere e scrivere i
\emph{parametri} del veicolo, eseguire le calibrazioni e monitorare la telemetria
in tempo reale.

Il collegamento tra QGC e il Pixhawk \`e sempre avvenuto attraverso il
\textbf{link radio di telemetria a 433\,MHz}, sia a banco sia in campo: la coppia
di moduli Holybro fa da ponte seriale wireless tra il computer e il flight
controller, e non \`e mai stato necessario collegare il drone via cavo USB. Questa
scelta ha il vantaggio di lavorare sempre nella stessa configurazione con cui il
drone vola, senza spostare cavi tra banco e campo.

Il capitolo ripercorre le configurazioni che hanno reso il drone operativo e
strumentato per l'acquisizione dati. Diamo spazio anche alle \textbf{limitazioni
e alle problematiche} incontrate: anche queste fanno parte a pieno titolo del
lavoro svolto, perch\'e documentano i limiti reali della piattaforma.

\section{Calibrazioni e setup iniziale}

Il \emph{commissioning} (la prima messa a punto del veicolo) parte dalla
selezione dell'\textbf{airframe}, cio\`e il modello geometrico con cui PX4 sa
quanti motori ci sono e come sono disposti: nel nostro caso la configurazione
esacottero \texttt{hexa\_x} (sei motori a X). Da qui PX4 ricava la
\emph{matrice di mixing}, ossia come tradurre i comandi di assetto in spinte sui
singoli motori.

Si procede poi con la \textbf{calibrazione dei sensori} --- accelerometri,
giroscopi e magnetometri --- e con la \textbf{calibrazione della radio}, che
associa le posizioni degli stick ai canali di comando. Infine si definiscono le
\emph{modalit\`a di volo} assegnandole agli interruttori del radiocomando.

Per le prove a banco, in cui non si vuole (o non si pu\`o) avere un aggancio GPS,
\`e stato abilitato l'\textbf{arming senza fix GPS}
(\texttt{COM\_ARM\_WO\_GPS\,=\,1}) in modalit\`a \emph{Stabilized}; il parametro
\`e stato riportato a \texttt{0} per i voli all'aperto, dove la posizione \`e
indispensabile e armare senza GPS sarebbe pericoloso.

\figurasegnaposto{Schermate QGroundControl delle fasi di setup: Airframe,
Sensors/Calibration, Flight Modes.}{QGC -- da catturare}{fig:qgc-config}

\section{Migrazione del GPS da CAN a UART}

Il modulo GNSS (u-blox NEO-M8P, integrato nel CubePilot Here+) era originariamente
interfacciato al flight controller \textbf{Cube Black} tramite bus
\textbf{DroneCAN}, un protocollo seriale a due fili condiviso tra pi\`u
dispositivi. La sostituzione del controllore con il Pixhawk 6X, resa necessaria dal
guasto del Cube Black (Capitolo~\ref{cap:diagnostici}), ha reso non pi\`u
utilizzabile quel collegamento: non si disponeva di un cablaggio CAN compatibile
con il nuovo controllore. Si \`e quindi optato per un collegamento su porta
\textbf{UART} dedicata (la GPS2), adattando un cavo con un connettore realizzato su
misura. La soluzione, oltre a rendere possibile l'integrazione, offre un vantaggio
diagnostico: l'UART \`e un collegamento punto-punto, privo dell'overhead del bus e
pi\`u agevole da analizzare a livello di flusso seriale. Lo stesso cavo \`e stato
successivamente oggetto di un intervento migliorativo, documentato nel medesimo
capitolo.

Subito dopo la migrazione il GPS non agganciava pi\`u alcun satellite: la console
di QGC riportava \texttt{baudrate: 0} e \texttt{rate reading: 0\,B/s}, segno che
sulla porta non transitava alcun dato. La causa era duplice --- un residuo della
configurazione DroneCAN ancora attiva e un disallineamento tra la \emph{porta
logica} dichiarata nei parametri e la \emph{porta fisica} a cui il cavo era
collegato. La soluzione \`e nei parametri di Tabella~\ref{tab:gps-migr}, seguiti
da un \emph{cold start} all'aperto (primo aggancio da fermo, a cielo libero).

\begin{table}[H]
\centering
\caption{Parametri risolutivi della migrazione GPS da DroneCAN a UART.}
\label{tab:gps-migr}
\small
\begin{tabularx}{\linewidth}{@{}l l L{7.2cm}@{}}
\toprule
\textbf{Parametro} & \textbf{Valore} & \textbf{Motivazione} \\
\midrule
\texttt{UAVCAN\_ENABLE} & 0 & Disabilita DroneCAN, in conflitto col driver UART \\
\texttt{GPS\_1\_CONFIG} & 202 (GPS\,2) & Allinea la porta logica al cavo fisico (era 201 = GPS\,1) \\
\texttt{GPS\_1\_PROTOCOL} & 1 (UBX) & Protocollo binario u-blox \\
\bottomrule
\end{tabularx}
\end{table}

\begin{notabox}[Nota di configurazione]
Ogni cambio di firmware o riconfigurazione della porta GPS richiede di
riverificare la corrispondenza tra il parametro \texttt{GPS\_x\_CONFIG} e la porta
fisica usata. Va inoltre segnalato che sul Pixhawk 6X non esistono i parametri
\texttt{CAN\_Px\_DRIVER} citati in parte della documentazione: per disattivare il
bus \`e sufficiente \texttt{UAVCAN\_ENABLE\,=\,0}.
\end{notabox}

\figurasegnaposto{Schema della migrazione GPS DroneCAN$\rightarrow$UART: cablaggio
prima e dopo.}{Schema -- da generare}{fig:gps-can-uart}

\section{Power module e calibrazione della batteria}

Il \emph{power module} \`e il sensore che misura tensione e corrente assorbite
dalla batteria e le pubblica al firmware. Inizialmente non pubblicava affatto il
topic \texttt{battery\_status} (lo stato risultava \emph{never published}): il
firmware non sapeva su quali ingressi analogici leggere i due valori. Assegnando
i canali ADC corretti (\texttt{BAT1\_V\_CHANNEL\,=\,16} per la tensione,
\texttt{BAT1\_I\_CHANNEL\,=\,17} per la corrente) e completando la calibrazione di
tensione, la telemetria di tensione e corrente \`e diventata disponibile a
$\sim$5\,Hz.

La taratura della batteria \`e tuttavia rimasta \textbf{problematica}. Il firmware
continua a segnalare alcuni \textbf{parametri mancanti} nella configurazione del
pacco e, di conseguenza, mostra in QGC un \textbf{livello di carica fisso}, non
legato allo stato reale della batteria. Lo stesso si riflette nei log: la
\textbf{percentuale di carica (SoC) registrata non \`e affidabile} e non
rispecchia la scarica effettiva. Per questo motivo, durante la campagna di volo,
l'autonomia residua \`e stata valutata in modo indiretto --- sulla
\textbf{tensione di pacco} e sul tempo di volo --- e non sulla percentuale
riportata dal firmware (vedi anche Capitolo~\ref{cap:campagna}). Tensione e
corrente istantanee restano comunque corrette e utilizzabili.

\section{Telemetria ESC su TELEM2}

Gli ESC (\emph{Electronic Speed Controller}, i regolatori che pilotano i motori)
sono Tekko32 F4 con firmware AM32 e sono comandati in \textbf{DShot}, un
protocollo digitale che invia i comandi di spinta senza le imprecisioni del
vecchio segnale PWM analogico. Ogni ESC restituisce inoltre la propria
\textbf{telemetria} (la cosiddetta telemetria KISS) su un filo dedicato; i sei
fili sono uniti su un'unica linea seriale (TELEM2). Per evitare collisioni sul bus
condiviso, il flight controller interroga gli ESC uno alla volta
(\emph{request-driven}): chiede il dato e attende la risposta, in sequenza.

La configurazione \`e riassunta in Tabella~\ref{tab:esc-tel}. Il risultato \`e il
topic \texttt{esc\_status} a $\sim$65\,Hz, con regime (RPM), tensione, corrente e
temperatura di ciascuno dei sei motori --- dati centrali per individuare squilibri
ed eliche danneggiate (Capitolo~\ref{cap:campagna}).

\begin{table}[H]
\centering
\caption{Parametri della telemetria ESC e del logging ad alta frequenza.}
\label{tab:esc-tel}
\small
\begin{tabularx}{\linewidth}{@{}l l L{7.2cm}@{}}
\toprule
\textbf{Parametro} & \textbf{Valore} & \textbf{Significato} \\
\midrule
\texttt{DSHOT\_CONFIG} & 600 & Protocollo DShot600 di comando agli ESC \\
\texttt{DSHOT\_TEL\_CFG} & 102 (TELEM2) & Porta seriale su cui leggere la telemetria KISS \\
\texttt{SDLOG\_PROFILE} & 849 & Profilo di logging (vedi Cap.~\ref{cap:acquisizione}) \\
\texttt{IMU\_GYRO\_RATEMAX} & 1000\,Hz & Frequenza massima del giroscopio \\
\bottomrule
\end{tabularx}
\end{table}

\section{Modalit\`a di volo e failsafe}

Nei test sono state usate due modalit\`a di volo. In \emph{Stabilized} il pilota
comanda direttamente l'assetto (inclinazione) del drone, mentre la quota \`e
gestita manualmente con il throttle; \`e la modalit\`a usata per gli hover e per
le prove con le pale. In \emph{Position} (POSCTL) il velivolo mantiene
attivamente la propria posizione usando il GPS, e a stick fermi resta sul punto;
\`e la modalit\`a usata per le traiettorie.

Un \textbf{failsafe} \`e una reazione automatica del firmware a una condizione
anomala. Quello pi\`u rilevante per noi \`e il failsafe di \textbf{perdita di
posizione}: se la stima di posizione diventa inaffidabile, PX4 pu\`o forzare un
atterraggio automatico \enquote{alla cieca} (\emph{blind-land}). A seguito degli
eventi descritti nel Capitolo~\ref{cap:diagnostici}, abbiamo rivisto questo
comportamento. Il punto chiave \`e che la decisione reale di blind-land nasce a
livello del filtro di stima \textbf{EKF} (il filtro che fonde IMU, GPS e
barometro per stimare assetto e posizione), tramite il parametro
\texttt{EKF2\_NOAID\_TOUT}, e non solo nei parametri del \emph{commander}: agire
solo su questi ultimi non basta. La Tabella~\ref{tab:failsafe} riassume la
configurazione adottata, orientata a concedere pi\`u tempo prima del failsafe e a
ridurre la deriva orizzontale durante un eventuale evento.

\begin{table}[H]
\centering
\caption{Configurazione del failsafe di perdita posizione e dei limiti di volo.}
\label{tab:failsafe}
\small
\begin{tabularx}{\linewidth}{@{}l c c L{6.2cm}@{}}
\toprule
\textbf{Parametro} & \textbf{Default} & \textbf{Applicato} & \textbf{Motivazione} \\
\midrule
\texttt{EKF2\_NOAID\_TOUT} & 5\,s & \textbf{10\,s} & Ritarda \texttt{xy\_valid=false}, vera causa del blind-land \\
\texttt{COM\_POSCTL\_NAVL} & Land & \textbf{Altitude} & Lascia gli stick orizzontali attivi al pilota \\
\texttt{COM\_POS\_FS\_EPH} & 5\,m & 10\,m & Soglia di accuratezza orizzontale del commander \\
\texttt{MPC\_XY\_VEL\_MAX} & 12\,m/s & \textbf{5\,m/s} & Riduce la deriva orizzontale massima \\
\texttt{MPC\_TILTMAX\_AIR} & 45\textdegree & \textbf{25\textdegree} & Limita inclinazione e velocit\`a massima \\
\bottomrule
\end{tabularx}
\end{table}

\section{Limitazioni e problematiche riscontrate}

La configurazione ha avuto alcuni limiti che merita documentare:

\begin{itemize}
  \item \textbf{RTK non utilizzabile}: la correzione RTK, che porterebbe la
        precisione di posizione a livello centimetrico, non \`e stata sfruttabile
        perch\'e il driver GPS va in crash alla ricezione dei \emph{fix}
        (Capitolo~\ref{cap:diagnostici}). La campagna \`e quindi stata condotta
        con il solo GPS standard ($\sim$80\,cm di precisione).
  \item \textbf{Parametri non disponibili nella build PX4}: alcuni parametri di
        failsafe documentati nei manuali (\texttt{COM\_POS\_FS\_DELAY},
        \texttt{COM\_POS\_FS\_EPV}) non esistono in questa versione del firmware;
        altri (\texttt{MPC\_JERK\_MAX}, \texttt{MPC\_ACC\_HOR}) non sono
        modificabili a mano perch\'e vincolati automaticamente dal \emph{trajectory
        shaping} (la generazione interna della traiettoria). In QGC gli slider
        \emph{Flight Behavior} agiscono come comandi di alto livello su questi
        gruppi di parametri.
  \item \textbf{Override RC spurio sul canale yaw}: in un volo automatico la
        missione \`e stata interrotta da un \emph{takeover} (passaggio forzato al
        controllo manuale) non voluto, innescato da un piccolo disturbo
        ($\sim$10\,\%) sul canale yaw a stick fermi. Pur sotto la soglia nominale
        \texttt{COM\_RC\_STICK\_OV}, il disturbo \`e bastato a far scattare
        l'override. Le mitigazioni possibili sono allargare la \emph{deadzone} del
        canale, alzare la soglia o disabilitare l'override in modalit\`a
        automatica.
\end{itemize}

\figurasegnaposto{Schermata QGC di una problematica riscontrata (es.\ qualit\`a
RTK insufficiente o messaggio di failsafe).}{QGC -- da catturare}{fig:qgc-problema}
